
\documentclass[11pt]{article}
\usepackage[margin=1in]{geometry}
\usepackage[T1]{fontenc}
\usepackage[utf8]{inputenc}
\usepackage{lmodern}
\usepackage{amsmath}
\usepackage{amssymb}
\usepackage{graphicx}
\usepackage{booktabs}
\usepackage{float}
\usepackage{hyperref}
\hypersetup{hidelinks}
\setlength{\parskip}{0.5em}
\setlength{\parindent}{0pt}
\begin{document}
\sloppy
\section{Step 5: Fixed-Window Onset-Locked ROI Follow-Up Analysis}
\label{sec:step5_fixed_window}
% Analyst note added 2026-03-21: a separate Step 5 visualization report was generated in parallel to the main Step 5 report. Because no additional example code snippet for the requested extra plot was present in the workspace or chat at implementation time, the extra supplementary plot was implemented as an Abstract-minus-Concrete HbO joint topography over time, using the same onset-locked epoching and plotting framework as the Step 4-style visualization set.
% Analyst note added 2026-03-21: the later expansion of the supplementary visualization report adapted example plots originally framed as Tapping/Control and Left/Right motor-task contrasts to the present experiment's Abstract versus Concrete conditions. To stay aligned with the Step 5 analysis rules, those descriptive plots were also restricted to the long-channel montage used in the fixed-window ROI analysis.

\subsection{Objective}
The objective of Step~5 is to perform a simple, targeted, onset-locked ROI analysis that is more
aligned with the time-localized fNIRS pattern suggested by the Step~4 visualizations.

This step is designed to answer the following question:
\begin{quote}
When the fNIRS response is summarized in a fixed post-onset window centered on the delayed HbO
response, do abstract and concrete questions differ in the pre-defined left frontal and left posterior
ROIs?
\end{quote}

Step~5 should produce:
\begin{itemize}
    \item a primary p-value for the front ROI;
    \item secondary p-values for the back ROI and the front-versus-back dissociation;
    \item confidence intervals and effect sizes;
    \item a short conclusion written in clear language.
\end{itemize}

\subsection{Status of This Step}
Step~5 is a \textbf{targeted follow-up analysis}.
It does \textbf{not} replace the pre-specified Step~4 primary test.
Instead, it follows from the observation that:
\begin{itemize}
    \item the Step~4 whole-question-window dissociation test was not statistically significant;
    \item the front ROI showed only a weak descriptive abstract-greater-than-concrete tendency;
    \item the onset-locked visualizations suggested that the condition difference is more localized
    in time than the whole-question-window metric.
\end{itemize}

Therefore, Step~5 should be described as:
\begin{quote}
a focused onset-locked follow-up designed to test whether the abstract-versus-concrete effect is
stronger in a fixed post-onset HbO window than in the whole-question-window average used in Step~4.
\end{quote}

\subsection{Why Step~5 Is Needed}
In Step~4, the trial metric averaged the signal from:
\begin{verbatim}
question_start --> answer
\end{verbatim}

This whole-question-window approach is reasonable, but it has one important limitation:
participants answered at different times, and the haemodynamic response is delayed.
As a result, a condition difference that is concentrated around a specific post-onset interval can be
diluted when averaging over a long and variable trial duration.

The onset-locked supplementary visualizations suggest exactly this possibility:
the condition difference appears more time-localized than sustained over the entire question window.
Step~5 is therefore a natural next step because it changes only one thing:
\begin{quote}
the summary metric is moved from a variable question-duration mean to a fixed onset-locked time window.
\end{quote}

\subsection{Core Analytical Decision}
To keep Step~5 simple, the following elements should remain unchanged from Step~4:
\begin{enumerate}
    \item the same included participant-sessions;
    \item the same preprocessed files from Step~3;
    \item the same long-channel-only rule;
    \item the same front and back ROI definitions;
    \item the same exclusion of short channels and transition channels from the primary ROI summaries;
    \item HbO as the primary chromophore.
\end{enumerate}

The main change in Step~5 is:
\begin{quote}
use a fixed post-onset HbO analysis window instead of the full question-start-to-answer interval.
\end{quote}

\subsection{Primary Scientific Question}
The primary Step~5 question is:
\begin{quote}
In the left frontal ROI, is the fixed-window onset-locked HbO response significantly different
between abstract and concrete questions?
\end{quote}

This is the most targeted and most plausible follow-up, because the Step~4 onset-locked plots suggest
that the strongest condition separation may be frontal and delayed after question onset, rather than
distributed uniformly across both front and back ROIs.

\subsection{Inputs}
Step~5 begins from:
\begin{itemize}
    \item the preprocessed participant-session fNIRS files from Step~3;
    \item the harmonized trigger table from Step~3;
    \item the inclusion table and ROI definitions from Step~4;
    \item the question metadata and abstract/concrete labels from Step~2.
\end{itemize}

\subsection{Frozen ROI Definitions}
Step~5 should use exactly the same ROI definitions as Step~4.

\subsubsection{Front ROI HbO channels}
\begin{verbatim}
S2_D6 hbo
S2_D4 hbo
S1_D6 hbo
S1_D3 hbo
S5_D6 hbo
S5_D4 hbo
S5_D3 hbo
\end{verbatim}

\subsubsection{Back ROI HbO channels}
\begin{verbatim}
S4_D2 hbo
S3_D1 hbo
S6_D2 hbo
S6_D1 hbo
S6_D5 hbo
S7_D2 hbo
S7_D5 hbo
S8_D1 hbo
S8_D5 hbo
\end{verbatim}

\subsubsection{Transition HbO channels excluded from the primary confirmatory summaries}
\begin{verbatim}
S5_D7 hbo
S4_D4 hbo
S4_D7 hbo
S3_D3 hbo
S3_D7 hbo
S6_D7 hbo
\end{verbatim}

\subsubsection{Short-separation pairs excluded from inference}
\begin{verbatim}
S1_D8
S2_D9
S3_D10
S4_D11
S5_D12
S6_D13
S7_D14
S8_D15
\end{verbatim}

These short pairs are used only for preprocessing and nuisance correction and must not be included
in the inferential ROI analysis.

\subsection{Epoching Strategy}
Step~5 should use the onset-locked epoch structure already shown to be feasible in the supplementary
Step~4 visualizations.
The epoch definition should therefore remain:

\begin{verbatim}
epoch window   = [-3.5 s, 13.2 s] around question onset
baseline       = [-2.0 s, 0.0 s]
event lock     = question_start
\end{verbatim}

This choice is advantageous because:
\begin{itemize}
    \item it already worked with zero dropped trials in the supplementary onset-locked analysis;
    \item it captures the delayed haemodynamic response;
    \item it keeps all trials in a common temporal frame.
\end{itemize}

\subsection{Primary Fixed Analysis Window}
The primary Step~5 HbO window should be:
\begin{verbatim}
[7.0 s, 11.0 s] after question onset
\end{verbatim}

This is the recommended primary window because:
\begin{itemize}
    \item it is centered on 9.0 s, which was already used in the supplementary single-time comparison;
    \item it falls within the delayed post-onset period where the visualizations suggest condition
    differences are more visible;
    \item it is wide enough to be stable, but still focused enough to avoid unnecessary dilution.
\end{itemize}

\subsection{Important Interpretation Rule}
Step~5 should be described honestly as a targeted follow-up.
Therefore:
\begin{itemize}
    \item the Step~4 result remains the pre-specified primary analysis result;
    \item the Step~5 fixed-window result is a post-Step~4 follow-up analysis;
    \item any significant finding in Step~5 should be described as follow-up evidence rather than
    a replacement of Step~4.
\end{itemize}

\subsection{Participant Inclusion Rules}
A participant-session should be included in Step~5 if:
\begin{itemize}
    \item it was included in Step~4;
    \item the onset trigger is available and valid;
    \item the full epoch window \([-3.5, 13.2]\) s is available;
    \item the baseline interval \([-2, 0]\) s is available;
    \item at least one valid front ROI HbO channel remains;
    \item at least one valid back ROI HbO channel remains.
\end{itemize}

Participants excluded from Step~4 should remain excluded unless a separate re-review explicitly
changes their status.

\subsection{Trial Inclusion Rules}
A trial should be included if:
\begin{itemize}
    \item it has a valid question onset event;
    \item it has a valid condition label (\texttt{abstract} or \texttt{concrete});
    \item the epoch window is fully available;
    \item the baseline interval is fully available;
    \item at least one valid ROI channel is present for the ROI being summarized.
\end{itemize}

Unlike Step~4, the primary Step~5 metric does \textbf{not} require the answer timestamp to define
the analysis window, because the step is specifically onset-locked rather than answer-locked.

\subsection{Step-by-Step Workflow}

\subsubsection{Step~5.1: Reuse the Step~4 inclusion set}
Start from the Step~4 included participant-session set.
Do not redefine the cohort unless a clear file or QC problem is discovered.

\subsubsection{Step~5.2: Construct onset-locked epochs}
For each valid trial:
\begin{itemize}
    \item lock the epoch to \texttt{question\_start};
    \item extract the time interval from \(-3.5\) s to \(13.2\) s;
    \item apply the baseline correction using \([-2.0, 0.0]\) s.
\end{itemize}

\subsubsection{Step~5.3: Compute trial-level fixed-window HbO values}
For each trial and each ROI, compute the mean baseline-corrected HbO value in:
\begin{verbatim}
[7.0 s, 11.0 s]
\end{verbatim}

\subsubsection{Step~5.4: Average channels within ROI}
For each trial:
\begin{itemize}
    \item average valid front ROI channels to obtain one front ROI trial value;
    \item average valid back ROI channels to obtain one back ROI trial value.
\end{itemize}

If a channel is missing or marked bad, omit it from the ROI average and record how many valid channels
contributed to the trial-level ROI summary.

\subsubsection{Step~5.5: Build participant-level condition means}
For each participant, compute:
\begin{itemize}
    \item mean front ROI value for abstract trials;
    \item mean front ROI value for concrete trials;
    \item mean back ROI value for abstract trials;
    \item mean back ROI value for concrete trials.
\end{itemize}

\subsubsection{Step~5.6: Run the primary front-ROI test}
Run one two-sided paired t-test across participants comparing:
\begin{equation}
\bar{y}^{(F,A)}_i \quad \text{versus} \quad \bar{y}^{(F,C)}_i
\end{equation}

This is the primary inferential test for Step~5.

\subsubsection{Step~5.7: Run the secondary back-ROI test}
Run one two-sided paired t-test across participants comparing:
\begin{equation}
\bar{y}^{(B,A)}_i \quad \text{versus} \quad \bar{y}^{(B,C)}_i
\end{equation}

This is a secondary test.

\subsubsection{Step~5.8: Run the secondary dissociation test}
For each participant, compute:
\begin{equation}
\Delta_i^{(7\text{--}11)}
=
\bigl(\bar{y}^{(F,A)}_i - \bar{y}^{(F,C)}_i\bigr)
-
\bigl(\bar{y}^{(B,A)}_i - \bar{y}^{(B,C)}_i\bigr)
\end{equation}

Then run a one-sample t-test on \(\Delta_i^{(7\text{--}11)}\) against zero.
This should be reported as a secondary dissociation analysis, not as the primary test.

\subsubsection{Step~5.9: Run simple sensitivity analyses}
To keep the workflow simple but still robust, include the following sensitivity checks:

\paragraph{Sensitivity window 1}
Repeat the analysis using:
\begin{verbatim}
[6.0 s, 10.0 s]
\end{verbatim}

\paragraph{Sensitivity window 2}
Repeat the analysis using:
\begin{verbatim}
[8.0 s, 12.0 s]
\end{verbatim}

\paragraph{Single-time-point sensitivity}
Repeat the ROI comparison at:
\begin{verbatim}
9.0 s
\end{verbatim}

\paragraph{Fast-response sensitivity}
Repeat the primary front ROI analysis after excluding trials whose behavioral answer occurred
before the end of the primary fixed window, that is:
\begin{verbatim}
answer time < 11.0 s after question onset
\end{verbatim}

This sensitivity analysis is useful because it assesses whether the fixed-window result depends
heavily on post-answer signal in very fast trials.

\subsection{Primary Outcome}
The primary Step~5 outcome is the participant-level front ROI fixed-window contrast:
\begin{equation}
d^{(F)}_i = \bar{y}^{(F,A)}_i - \bar{y}^{(F,C)}_i
\end{equation}

The primary Step~5 hypothesis is:
\begin{quote}
the front ROI fixed-window HbO response differs between abstract and concrete questions.
\end{quote}

\subsection{Trial-Level Signal Definition}
For participant \(i\), channel \(j\), and trial \(k\), let \(x_{ijk}(t)\) denote the preprocessed
HbO signal in the onset-locked epoch after baseline correction.
Then the fixed-window trial summary is:
\begin{equation}
r^{(7\text{--}11)}_{ijk}
=
\overline{x_{ijk}(t \in [7,11])}
\end{equation}

Since the epoch is already baseline-corrected using \([-2,0]\) s, no additional subtraction is needed here.

\subsection{ROI-Level Trial Summaries}
For each trial, define:

\paragraph{Front ROI trial value}
\begin{equation}
y^{(F)}_{ik}
=
\frac{1}{J^{(F)}_{ik}}
\sum_{j \in F_{ik}} r^{(7\text{--}11)}_{ijk}
\end{equation}

\paragraph{Back ROI trial value}
\begin{equation}
y^{(B)}_{ik}
=
\frac{1}{J^{(B)}_{ik}}
\sum_{j \in B_{ik}} r^{(7\text{--}11)}_{ijk}
\end{equation}

where \(F_{ik}\) and \(B_{ik}\) are the sets of valid front and back ROI channels contributing
to trial \(k\) for participant \(i\).

\subsection{Participant-Level Condition Means}
For each participant:
\begin{equation}
\bar{y}^{(F,A)}_i
=
\frac{1}{n_i^{(A)}}
\sum_{k \in A} y^{(F)}_{ik},
\qquad
\bar{y}^{(F,C)}_i
=
\frac{1}{n_i^{(C)}}
\sum_{k \in C} y^{(F)}_{ik}
\end{equation}

\begin{equation}
\bar{y}^{(B,A)}_i
=
\frac{1}{n_i^{(A)}}
\sum_{k \in A} y^{(B)}_{ik},
\qquad
\bar{y}^{(B,C)}_i
=
\frac{1}{n_i^{(C)}}
\sum_{k \in C} y^{(B)}_{ik}
\end{equation}

\subsection{Primary Statistical Test}
The primary Step~5 test should be a \textbf{two-sided paired t-test} across participants on:
\begin{equation}
d^{(F)}_i = \bar{y}^{(F,A)}_i - \bar{y}^{(F,C)}_i
\end{equation}

\paragraph{Null hypothesis}
\begin{equation}
H_0: \mu_{d^{(F)}} = 0
\end{equation}

\paragraph{Alternative hypothesis}
\begin{equation}
H_1: \mu_{d^{(F)}} \neq 0
\end{equation}

The significance threshold should be:
\begin{verbatim}
alpha = 0.05
\end{verbatim}

\subsection{Why the Front ROI Is Primary in Step~5}
The front ROI is the most appropriate primary target in this follow-up step because:
\begin{itemize}
    \item the original project hypothesis predicts stronger abstract-related activation in prefrontal regions;
    \item the Step~4 whole-window analysis showed a weak descriptive front tendency but not a significant result;
    \item the onset-locked waveforms suggest that any condition difference may be delayed and more visible
    in a fixed post-onset window.
\end{itemize}

\subsection{Secondary Tests}
Step~5 should report the following as secondary:
\begin{enumerate}
    \item paired t-test for the back ROI;
    \item one-sample t-test for the front-versus-back dissociation score;
    \item the same three tests for HbR;
    \item the sensitivity windows and the single-time-point 9.0 s check.
\end{enumerate}

\subsection{Primary Reporting Quantities}
The main Step~5 results table must report:
\begin{itemize}
    \item number of included participants;
    \item number of valid abstract and concrete trials used;
    \item mean front ROI HbO value for abstract trials;
    \item mean front ROI HbO value for concrete trials;
    \item mean front ROI difference;
    \item paired t statistic;
    \item degrees of freedom;
    \item exact p-value;
    \item 95\% confidence interval for the mean front ROI difference;
    \item Cohen's \(d_z\) for the paired contrast.
\end{itemize}

\subsection{Secondary Reporting Quantities}
The secondary results table should report:
\begin{itemize}
    \item back ROI abstract and concrete means;
    \item back ROI paired t-test result;
    \item dissociation score mean and one-sample t-test result;
    \item HbR analogues;
    \item sensitivity-window results;
    \item the 9.0 s single-time-point results.
\end{itemize}

\subsection{Primary Conclusion Rule}
The final Step~5 conclusion must be based first on the front ROI fixed-window HbO test.

The rule should be:
\begin{itemize}
    \item if \(p < 0.05\) and the front ROI mean difference is positive, conclude that abstract
    questions elicited a significantly larger onset-locked HbO response than concrete questions
    in the left frontal ROI during the fixed post-onset window;
    \item if \(p < 0.05\) and the front ROI mean difference is negative, conclude that concrete
    questions elicited a significantly larger onset-locked HbO response than abstract questions
    in the left frontal ROI during the fixed post-onset window;
    \item if \(p \geq 0.05\), conclude that the fixed-window onset-locked front ROI analysis
    did not show a statistically significant abstract-versus-concrete difference.
\end{itemize}

\subsection{Conclusion Templates}

\paragraph{Template if significant and abstract \(>\) concrete in the front ROI.}
In the targeted fixed-window onset-locked follow-up analysis, abstract questions elicited a
significantly larger HbO response than concrete questions in the left frontal ROI during the
7--11~s post-onset window
(\(t(df)=\ldots\), \(p=\ldots\), mean difference \(=\ldots\), 95\% CI \([\ldots,\ldots]\), \(d_z=\ldots\)).
This suggests that the abstract-versus-concrete effect is more visible in a delayed onset-locked
analysis than in the whole-question-window metric used in Step~4.

\paragraph{Template if significant and concrete \(>\) abstract in the front ROI.}
In the targeted fixed-window onset-locked follow-up analysis, concrete questions elicited a
significantly larger HbO response than abstract questions in the left frontal ROI during the
7--11~s post-onset window
(\(t(df)=\ldots\), \(p=\ldots\), mean difference \(=\ldots\), 95\% CI \([\ldots,\ldots]\), \(d_z=\ldots\)).

\paragraph{Template if not significant.}
In the targeted fixed-window onset-locked follow-up analysis, the difference between abstract and
concrete questions in the left frontal ROI during the 7--11~s post-onset window was not
statistically significant
(\(t(df)=\ldots\), \(p=\ldots\), mean difference \(=\ldots\), 95\% CI \([\ldots,\ldots]\), \(d_z=\ldots\)).

\subsection{Expected Outputs}

\subsubsection*{Tables}
\begin{verbatim}
01_step5_inclusion_table.csv
02_step5_epoch_manifest.csv
03_step5_trial_roi_summary_hbo.csv
04_step5_trial_roi_summary_hbr.csv
05_step5_participant_roi_summary.csv
06_step5_primary_front_test.csv
07_step5_secondary_back_test.csv
08_step5_secondary_dissociation_test.csv
09_step5_sensitivity_windows.csv
10_step5_fast_response_sensitivity.csv
11_step5_final_conclusion.csv
\end{verbatim}

\subsubsection*{Figures}
\begin{verbatim}
figures/step5/front_roi_window_plot.png
figures/step5/back_roi_window_plot.png
figures/step5/front_waveform_with_window.png
figures/step5/back_waveform_with_window.png
figures/step5/dissociation_histogram_fixed_window.png
figures/step5/topography_9s.png
\end{verbatim}

\subsubsection*{Reports}
\begin{verbatim}
reports/step5/step5_fixed_window_report.tex
reports/step5/tables/primary_front_test_table.tex
reports/step5/tables/secondary_tests_table.tex
reports/step5/text/final_conclusion.tex
\end{verbatim}

\subsection{Minimal Figures for the Main Report}
The main Step~5 report should include at least:
\begin{enumerate}
    \item a paired participant plot for the front ROI fixed-window contrast;
    \item a paired participant plot for the back ROI fixed-window contrast;
    \item a waveform plot for the front ROI with the 7--11~s analysis window shaded;
    \item a waveform plot for the back ROI with the 7--11~s analysis window shaded;
    \item a histogram of the dissociation scores;
    \item a 9.0~s topographic comparison figure.
\end{enumerate}

\subsection{Quality-Control Checks}
Before Step~5 is considered complete, verify that:
\begin{enumerate}
    \item the same Step~4 included participant-sessions are used unless explicitly re-reviewed;
    \item the epoch window \([-3.5, 13.2]\) s is fully available for all included trials;
    \item the baseline window \([-2, 0]\) s is fully available;
    \item only long-separation channels are used in the inferential tests;
    \item the short-separation pairs are excluded from the ROI summaries;
    \item the transition channels are excluded from the primary ROI summaries;
    \item the front ROI channel list is identical across participants;
    \item the back ROI channel list is identical across participants;
    \item the primary 7--11~s front ROI test is run once and clearly labeled as primary;
    \item all other tests are labeled secondary or sensitivity analyses;
    \item exact p-values appear in the final report.
\end{enumerate}

\subsection{Minimal Completion Criteria}
Step~5 is complete when:
\begin{itemize}
    \item the onset-locked epochs have been created;
    \item the 7--11~s ROI summaries have been computed;
    \item the participant-level front and back ROI condition means exist;
    \item the primary front ROI fixed-window test has been run;
    \item the secondary back and dissociation tests have been run;
    \item the sensitivity analyses have been completed;
    \item a short final conclusion has been generated;
    \item all exclusions and warnings are documented.
\end{itemize}

\subsection{Short Practical Summary}
In simple terms, Step~5 should:
\begin{quote}
keep the same participants, the same preprocessed data, and the same ROI definitions as Step~4,
but replace the whole-question-window metric with a fixed onset-locked 7--11~s HbO window,
test abstract versus concrete first in the front ROI, then in the back ROI and in the dissociation score,
report the p-values clearly, and state the conclusion honestly as a follow-up analysis.
\end{quote}

\end{document}
